\documentclass[12pt, a4paper]{article}

% --- IMPOSTAZIONI DI BASE E LINGUA ---
\usepackage[utf8]{inputenc}
\usepackage[T1]{fontenc}
\usepackage[italian]{babel}
\usepackage[none]{hyphenat}
\usepackage[a4paper, top=2.4cm, bottom=2.4cm, left=2cm, right=2cm, headheight=14pt, headsep=10pt]{geometry}

% --- TIPOGRAFIA RAFFINATA ---
\usepackage{microtype}
\setlength{\parindent}{0pt}
\setlength{\parskip}{6pt plus 1.5pt minus 1pt}
\setlength{\emergencystretch}{2em}

% --- MATEMATICA E FORMATTAZIONE ---
\usepackage{amsmath}
\usepackage{amssymb} % Aggiunto per i simboli come i numeri reali \mathbb{R}

% --- IMMAGINI, GRAFICA E COLORI ---
\usepackage{graphicx}
\usepackage{float}
\usepackage{xcolor} % Caricato PRIMA di tikz

% --- DISEGNO GRAFICI E POLIEDRI (TIKZ / PGFPLOTS) ---
\usepackage{tikz}
\usetikzlibrary{arrows}
\usepackage{pgfplots}
\pgfplotsset{compat=1.15} % Usiamo la 1.15 per massima stabilità
\usepgfplotslibrary{fillbetween}

% --- COLORI PERSONALIZZATI (Stile GeoGebra) ---
\definecolor{myblue}{rgb}{0,0,1}
\definecolor{myred}{rgb}{1,0,0}
\definecolor{mygreen}{rgb}{0,0.5,0}

% --- COLORE DI STILE (titoli, righe e box) ---
\definecolor{titlecolor}{RGB}{16,52,96}

% --- TITOLI DI SEZIONE ELEGANTI ---
\usepackage{titlesec}
\titleformat{\section}
  {\color{titlecolor}\Large\bfseries\raggedright}
  {\thesection}{0.7em}{}
  [\vspace{0.4ex}{\color{titlecolor!50}\titlerule[1pt]}]
\titleformat{\subsection}
  {\color{titlecolor}\large\bfseries\raggedright}
  {\thesubsection}{0.6em}{}
  []
\titlespacing*{\section}{0pt}{2.2em}{1em}
\titlespacing*{\subsection}{0pt}{1.8em}{0.6em}

% --- LISTE CON SPAZIATURE PIÙ ARIOSE ---
\usepackage{enumitem}
\setlist{itemsep=2pt, topsep=3pt, parsep=2pt, partopsep=0pt}

% --- TABELLE PIÙ RESPIRABILI ---
\renewcommand{\arraystretch}{1.25}

% --- DIDASCALIE RAFFINATE ---
\usepackage{caption}
\captionsetup{font=small, labelfont={bf,color=titlecolor}, labelsep=period, textfont=it, skip=4pt}

% --- TESTATINA E PIÈ DI PAGINA ---
\usepackage{zref-totpages}
\usepackage{fancyhdr}
\fancypagestyle{plain}{%
  \fancyhf{}%
  \fancyhead[L]{\small\itshape\color{gray}Ricerca Operativa}%
  \fancyhead[R]{\small\color{gray}Andrea Salvatore}%
  \fancyfoot[C]{\small\color{gray}Pagina \thepage\ di \ztotpages}%
  \renewcommand{\headrulewidth}{0.4pt}%
}
\pagestyle{plain}
\renewcommand{\headrulewidth}{0.4pt}
\renewcommand{\footrulewidth}{0pt}

% --- TITOLO DEL DOCUMENTO RAFFINATO ---
\makeatletter
\renewcommand{\@maketitle}{%
  \begin{center}%
    {\LARGE\bfseries\color{titlecolor}\@title\par}%
    \vskip 0.6em%
    {\large\color{titlecolor}\@author\par}%
    \vskip 0.5em%
    {\normalsize\itshape\color{gray}\@date\par}%
  \end{center}%
  \vskip 0.4em%
  {\color{titlecolor!50}\hrule height 1pt}%
  \vskip 1.2em%
}
\makeatother

% --- RIGA SEPARATRICE ELEGANTE ---
\newcommand{\separatore}{\noindent{\color{titlecolor!50}\rule{\textwidth}{0.6pt}}}

% --- BOX PER L'AVVISO DI OGNI LEZIONE ---
\newcommand{\avviso}[1]{%
  \par\medskip
  \noindent
  \begin{tikzpicture}
    \node[rounded corners=2.5pt, fill=titlecolor!5, draw=titlecolor!35, line width=0.6pt,
          inner xsep=9pt, inner ysep=7pt, text width=\dimexpr\linewidth-18.6pt\relax, align=left] {#1};
  \end{tikzpicture}
  \par\medskip
}

% --- LINK IPERTESTUALI ---
\usepackage[colorlinks=true, linkcolor=titlecolor, urlcolor=titlecolor, pdfauthor={Andrea Salvatore}, pdftitle={Ricerca Operativa}]{hyperref} % hyperref va sempre caricato per ultimo!

\begin{document}
	\title{\textbf{Ricerca Operativa}}
	\author{\textbf{Andrea Salvatore}}
	\date{Anno accademico: 2026/2027}
	\maketitle
	
	I miei contatti:
	\begin{itemize}
		\item \textbf{GitHub}: \href{https://github.com/andrea-salvatore}{andrea-salvatore}
		\item \textbf{Website}: \href{https://andreasalvatore.dev/}{andreasalvatore.dev}
		\item \textbf{Cellulare}: +39 3703116443
		\item \textbf{E-mail}:
		\begin{itemize}
			\item Personale: asalvatore.junior@gmail.com
			\item Lavorativa: ajs.sdk.mt@gmail.com
			\item Istituzionale: a.salvatore9@studenti.unipi.it
		\end{itemize}
	\end{itemize}
	\vspace{1cm}
	
	L'intento di questo file è contenere le lezioni del prof. Massimo Pappalardo in ordine strettamente cronologico seguendo, dunque, le lezioni dell'anno accademico citato sopra. L'idea nasce dal fatto che spesso capita che gli assenti abbiano bisogno di recuperare delle lezioni. Attraverso i file in questo repository ora hanno tutto in tempo reale lezione dopo lezione. Inoltre, per senso di collaborazione, agli studenti è fornito non solo il PDF, bensì anche il sorgente .tex, utile per correggere errori e/o aggiungere dettagli utili direttamente tramite le Pull Request di GitHub. Oppure, nell'ipotesi che io sia in difficoltà con la trascrizione corretta degli appunti, ci si può aiutare inserendo tutto qui e mettendo a disposizione di \textbf{tutti} il materiale in ordine cronologico e in tempi ristretti. Grazie a tutti per la collaborazione.
	
	\separatore

	\section*{Contenuto del corso}
	\begin{itemize}
		\item Programmazione Lineare (PL)
		\begin{itemize}
			\item Requisito: Algebra Lineare
		\end{itemize}
		\item Programmazione Lineare Intera (PLI)
		\begin{itemize}
			\item Requisito: Algebra Lineare
		\end{itemize}
		\item Programmazione Lineare su Reti (PL su Reti)
		\begin{itemize}
			\item Requisito: Algebra Lineare
		\end{itemize}
		\item Programmazione Non Lineare (PNL)
		\begin{itemize}
			\item Requisito: Analisi Matematica 2
		\end{itemize}
	\end{itemize}
	
	Libro di testo consigliato: \textbf{\textit{Ricerca Operativa}}\\
	Autori: \underline{Massimo Pappalardo} e \underline{Mauro Passacantando}
	\begin{figure}[H]
		\centering
		\includegraphics[scale = 1]{/home/ajs/Immagini/ro_libro.jpg}
		\caption{Libro di Ricerca Operativa del prof. Massimo Pappalardo}
		\label{fig:mio_grafico}
	\end{figure}
	
	\separatore\\
	
	\avviso{\textit{\textbf{Questa lezione, come le sue precedenti e successive, può contenere errori e/o imprecisioni. Per il bene di tutti, è buona norma far notare all'autore l'errore o l'imprecisione così da risolvere nel minor tempo possibile. I canali di comunicazione, per eventuali avvisi generali o di errori trovati negli appunti, sono in cima al file. Grazie per la collaborazione e Buono studio!}}}
		
	\section*{Lezione 1 | 21 settembre 2026} % Usato section* invece di andare a capo a mano
	\subsection*{Introduzione alla PL | Problema di produzione} % Usato subsection* invece di \title
	
	\textbf{Definizione di \underline{Problema di Programmazione Lineare}}:\\
	\textit{Si definisce tale un problema di minimo e massimo di una funzione lineare soggetta a vincoli lineari di $\leq$, $\geq$ o $=$. Un problema è Lineare se i suoi vincoli e la Funzione Obiettivo sono di primo grado}.\\
	\underline{Attenzione!!!} Scopriremo nella \textit{Lezione 2} che il $<$ e il $>$ non sono accettati perché la soluzione ottima del problema di PL è sempre sul bordo del Poliedro.\\
	
	Ipotizziamo un problema di questo tipo\\\\
	\begin{center}
		\begin{tabular}{|c|c|c|c|}
			\hline
			\textbf{Kg/Unità} & \textbf{Travi Reticolari} & \textbf{Pilastri Composti} & \textbf{Disponibilità} \\
			\hline
			\textbf{Acciaio} & 20 & 10 & 950\\
			\hline
			\textbf{Lavorazione} & 2 & 3 & 120\\
			\hline
			\textbf{Collaudo e verniciatura} & 1 & 2 & 70\\
			\hline
		\end{tabular}
	\end{center}
	Questa tabella indica le colonne come il bene da produrre (con le disponibilità finali) e le righe come le risorse da usare.\\\\
	Questo si chiama, appunto, "Problema di produzione" perché riguarda la realizzazione di un modello matematico applicato a lavori di produzione che hanno determinati vincoli. I medesimi vincoli si rappresentano attraverso un sistema che per la PL si chiama \textit{Formato Primale Standard} e si rappresenta così:
	\[
	\begin{aligned}
		& \begin{cases}
			\min/\max c^Tx\\
			Ax \leq b\\
			x \geq 0
		\end{cases}
	\end{aligned}
	\]
	\begin{itemize}
		\item $\boldsymbol{\min/\max c^Tx}$ $\rightarrow$ è la Funzione Obiettivo (FO) che ci indica cosa vogliamo minimizzare o massimizzare. È compreso in $R^n$ perché contiene solo il numero $n$ delle variabili del problema.
		\item $\boldsymbol{Ax \leq b}$ $\downarrow$
		\begin{itemize}
			\item \textbf{A} $\rightarrow$ è la Matrice che contiene i coefficienti del problema ed è compresa in $R^{m*n}$ con \textit{m} che indica il numero di vincoli del nostro problema e \textit{n} che indica il numero di variabili.
			\item \textbf{b} $\rightarrow$ è il vettore colonna che contiene i valori di $\leq$ associati rispettivamente al proprio vincolo. È compreso in $R^m$ perché contiene solo i valori a destra del $\leq$.
		\end{itemize}
		\item $\boldsymbol{x \geq 0}$ $\rightarrow$ Vincolo che ci indica il valore minimo di x dato che, in una situazione reale, non potremmo avere -10 travi reticolari, ad esempio. Serve per avere risultati realistici dato che il modello matematico considera \underline{SOLO} la matematica e non la realtà delle cose.
		\begin{itemize}
			\item $\boldsymbol{x_1}$ $\rightarrow$ Travi reticolari.
			\item $\boldsymbol{x_2}$ $\rightarrow$ Pilastri composti.
		\end{itemize}
	\end{itemize}
	Pappalardo ribadisce quanto sia importante la rappresentazione dei valori da dare al MATLAB per poter svolgere i calcoli (e farli svolgere al sistema) nella maniera corretta. Ad esempio, se avessimo 7 vincoli e 3 variabili, il numero totale di valori da dare al sistema sarà 31 perché:\\
	$\boldsymbol{(7 * 3) + 3 + 7 = 31}$\\
	Questi sono tutti i valori che fornirò al MATLAB per trovare l'ottimo, in particolare quei numeri si riferiscono a:
	\begin{itemize}
		\item $(7*3) \rightarrow$ che sarebbe la Matrice $A$.
		\item $+ 3 \rightarrow$ sommo il numero di variabili della 
		FO.
		\item $+ 7 \rightarrow$ sommo il numero dei risultati di ogni vincolo. Se abbiamo 7 vincoli di $\leq$, allora avrò 7 valori a destra del $\leq$, che quindi rappresentano il valore del vettore colonna $b$
	\end{itemize}
	
	L'esempio che possiamo fare riguarda l'esercizio della tabella precedente $\downarrow$
	
	\[ 
	\begin{aligned}
		\text{Problema in 2 Variabili $\rightarrow$ ($R^2$)}
		& \begin{cases}
			\max 300x_1 + 200x_2\\
			20x_1 + 10x_2 \leq 950\\
			2x_1 + 3x_2 \leq 120\\
			x_1 + 2x_2 \leq 70\\
			x_{1/2} \geq 0
		\end{cases}
	\end{aligned}
	\]
	
	Questa tipologia di problemi di ricerca operativa si definiscono come:
	
	\begin{itemize}
		\item \textbf{Ottimizzazione}: Nel caso del problema citato sopra, parliamo dell'ottimizzazione dei \textit{Guadagni} e la minimizzazione dei \textit{Costi}.
		\item \textbf{Risorse limitate}: Vincoli del problema. \textit{$Lineare = Proporzionale$}.
		\item \textbf{Variabili decisionali}: Trasformare il problema concreto in un modello matematico come mostrato prima.
		\begin{itemize}
			\item A seconda del problema si può intuire la correttezza del risultato ottenuto dal modello realizzato. Se il problema trattasse la produzione di indumenti come cappotti, pantaloni, e quant'altro, sappiamo che una soluzione ottima come $47,5$ non potrebbe mai essere corretta dato che non puoi produrre mezzo pantalone o mezzo cappotto, dunque, è corretto accettare valori come $47$ dato che è il valore intero più grande vicino alla soluzione ottima, ma se invece si sta parlando di liquidi, o in generale grammi o chilogrammi, allora il valore con la virgola esprime una quantità generica (e non un'unità come un pantalone appunto) e quindi $47,5$ sarebbe accettato.
		\end{itemize}
	\end{itemize}
	
	\underline{Piccola nota}: l'ultima riga dei vincoli del modello scritto qui su, quindi, "$x_{1/2} \geq 0$", in realtà rappresenta un modo veloce di scrivere:
	
	 \[
	 \begin{aligned}
	 	\begin{cases}
	 		x_1 + 0x_2 \geq 0\\
	 		0x_1 + x_2 \geq 0
	 	\end{cases}
	 \end{aligned}
	 \]
	Che è invece il modo corretto, considerando che poi le x diventeranno negative per cambiare il segno da $\geq$ a $\leq$, in questo modo diventa:
	\[
	\begin{aligned}
		\begin{cases}
			-x_1 + 0x_2 \leq 0\\
			0x_1 -x_2 \leq 0
		\end{cases}
	\end{aligned}
	\]
	Reputo importante saperlo perché questi valori fanno parte della matrice A spiegata in precedenza.
	\vspace{0.5cm}
	
	Attraverso il modello matematico che realizziamo, otteniamo dei risultati che, a seconda del valore stesso di questi risultati, possono essere:
	
	\begin{itemize}
		\item 	\textbf{Soluzione Ammissibile}: soluzione che rispetta tutti i vincoli del modello.
		\item \textbf{Soluzione Ottima}: La soluzione migliore possibile massima/minima che rispetta i vincoli.
	\end{itemize}
	
	\vspace{1cm}
	
	\textbf{Esercizio 1}\\
	\[
	\begin{aligned}
		\textbf{Problema in $\boldsymbol{R^3}$}
		\begin{cases}
			\max x_1 - 2x_2 + 3x_3\\
			6x_1 - 7x_2 \leq 9\\
			x_1 + 4x_2 + 5x_3 \geq 8
		\end{cases}
	\end{aligned}
	\]
	Domanda: è nella forma Primale Standard? No perché il secondo vincolo è di $\geq$ e va trasformato in $leq$ invertendo la funzione. Dunque, per portarlo in formato Primale Standard, devo:
	\[
	\begin{aligned}
		\begin{cases}
			\max x_1 - 2x_2 + 3x_3\\
			6x_1 - 7x_2 \leq 9\\
			\boldsymbol{-x_1 - 4x_2 - 5x_3 \leq -8}
		\end{cases}
	\end{aligned}
	\]
	
	\vspace{1cm}
	
	La funzione di MATLAB che usiamo per i problemi di PL (e poi anche di PLI e PL su Reti quando arriverà il loro monento) è \textbf{\underline{limprog}} che accetta valori in questo modo:\\\\ \textit{\textbf{linprog(c, A, b, Aeq, beq, lb, ub)}}\\\\
	dove \textit{Aeq} e \textit{beq} valgono 0 per i problemi di PL, dato che si usano in PLI, e poi, come detto in precedenza, abbiamo \textit{lb} che rappresenta il minimo valore che le x possono assumere e \textit{ub} che rappresenta il massimo valore che le x possono assumere. Le altre variabili non citate qui sono state spiegate in precedenza.
	
	\vspace{1cm}
	
	\textbf{IMPORTANTE!!!} \textit{linprog} cerca solo i massimi. Matematicamente sappiamo però che il massimo di una funzione è come fare il minimo della sua opposta:
	
	\textit{$\boldsymbol{f(x) = - f(x) \rightarrow}$}
	\underline{Ed è fondamentale da sapere!}\\
	Questo ci permette, dunque, di dire che moltiplicando le funzioni per -1 otteniamo:
	\begin{itemize}
		\item $\min \Leftrightarrow \max$
		\item $\leq \Leftrightarrow \geq$
	\end{itemize}
	Fine \textit{Lezione 1}
	
	\separatore\\
	
	\avviso{\textit{\textbf{Questa lezione, come le sue precedenti e successive, può contenere errori e/o imprecisioni. Per il bene di tutti, è buona norma far notare all'autore l'errore o l'imprecisione così da risolvere nel minor tempo possibile. I canali di comunicazione, per eventuali avvisi generali o di errori trovati negli appunti, sono in cima al file. Grazie per la collaborazione e Buono studio!}}}
	
	\section*{Lezione 2 | 22 settembre 2026}
	\subsection*{Continuazione del problema di produzione e introduzione del problema di assegnamento}
	
	Riprendendo da dove abbiamo lasciato: Se in un problema di PL ci trovassimo un vincolo con $=$, come si gestisce? Domanda che nasce sempre dal fatto che \textit{linprog} risolve solo problemi col $\leq$.\\
	Matematicamente, una funzione con $=$ si può rappresentare come la stessa funzione scritta due volte dove però la prima volta la poni con $\leq$ e la seconda con il $\geq$ che diventa nuovamente di $\leq$ sempre per le proprietà che abbiamo già spiegato.
	Un banale esempio potrebbe essere la funzione $\boldsymbol{5x_1 + 4x_2 = 6}$:
	
	\[
	\begin{aligned}
		\begin{cases}
			5x_1 + 4x_2 \leq 6\\
			5x_1 + 4x_2 \geq 6
		\end{cases}
	\end{aligned}
	\]
	e quindi...
	\[
	\begin{aligned}
		\begin{cases}
			5x_1 + 4x_2 \leq 6\\
			-5x_1 - 4x_2 \leq 6
		\end{cases}
	\end{aligned}
	\]
	Risolvendo così il problema dei vincoli di uguale ($=$) nei problemi di PL.\\
	
	In seguito vedremo il caso in cui si passa da $\leq$ a $=$ introducendo la variabile di \textit{Slack} (o Scarto).\\
	\textit{$\boldsymbol{f(x) \leq k \rightarrow f(x) + s = k}$}\\
	Ma si vedrà in seguito.
	
	\vspace{1cm}
		
	La risoluzione geometrica di un problema di PL richiede la conoscenza di questi 3 punti facenti parte del corso di \underline{Matematica 0}:
	\begin{itemize}
		\item Un $=$ nel piano rappresenta una Retta.
		\item Un $\leq$ nel piano rappresenta un semipiano inferiore.
		\item Un $\geq$ nel piano rappresenta un semipiano superiore.
	\end{itemize}
	Ricordiamo che all'esame si può usare tutto affinché \textbf{non} si usi la connessione internet, dunque, è possibile usare anche \textbf{\textit{GeoGebra}} che nel caso di questo corso è estremamente utile per la visualizzazione grafica del nostro modello matematico.\\
	
	Geometricamente, il grafico del problema affrontato precedentemente è questo:
	
	\begin{figure}[H]
		\centering
		\begin{tikzpicture}[line cap=round,line join=round]
			\begin{axis}[
				width=12cm,
				height=12cm,
				axis lines=middle,
				ymajorgrids=true,
				xmajorgrids=true,
				% ZOOM IMPOSTATO QUI
				xmin=-50,
				xmax=200,
				ymin=-50,
				ymax=200,
				xtick={0, 50, 100, 150},
				ytick={0, 50, 100, 150},
				xlabel = {$x_1$},
				ylabel = {$x_2$},
				]
				
				% Riquadro di clip allineato allo zoom
				\clip(-50,-50) rectangle (200,200);
				
				% VINCOLO 1 (Blu)
				\draw[line width=2pt, color=myblue, fill=myblue, fill opacity=0.15] 
				(-155.24,402.33)--(-396.58,402.33)--(-396.58,-427.61)--(262.88,-427.61)--(-155.24,402.33);
				
				% VINCOLO 2 (Rosso)
				\draw[line width=2pt, color=myred, fill=myred, fill opacity=0.15] 
				(-399.74,306.49)--(-396.58,-427.61)--(706.14,-427.61)--(-399.74,306.49);
				
				% VINCOLO 3 (Verde)
				\draw[line width=2pt, color=mygreen, fill=mygreen, fill opacity=0.15] 
				(-399.74,234.87)--(-396.58,-427.61)--(931.53,-427.61)--(-399.74,234.87);
				
				% VINCOLO DI NON NEGATIVITÀ x1 >= 0 (Grigio chiaro)
				\draw[line width=2pt, color=gray, fill=gray, fill opacity=0.1] 
				(0.09,402.33)--(0.09,-427.61)--(979.50,-427.61)--(979.50,402.33);
				
				% VINCOLO DI NON NEGATIVITÀ x2 >= 0 (Grigio scuro)
				\draw[line width=2pt, color=darkgray, fill=darkgray, fill opacity=0.1] 
				(-402.89,402.33)--(-402.89,-0.02)--(979.50,-0.02)--(979.50,402.33);
				
			\end{axis}
		\end{tikzpicture}
		\caption{La Regione Ammissibile è la zona dove tutti i colori si sovrappongono.}
	\end{figure}
	
	La Regione Ammissibile è anche detta \textbf{\textit{Poliedro}} la cui definizione è:
	\begin{itemize}
		\item Definizione che piace ai Matematici: \textit{Si definisce tale l'intersezione di un numero finito di semispazi chiusi.}
		\item Definizione che piace a Ingegneri e Informatici: \textit{Si definisce tale le soluzioni di $Ax \leq b$}.
	\end{itemize}
	Meglio usare quella che piace ai matematici dato che Pappalardo è un Matematico.\\\\
	\textbf{Nota importante sulla Notazione (per evitare disastri all'esame):}\\
	Il professore ci tiene a fare una precisazione fondamentale per non confondere i concetti. Fate molta attenzione all'uso delle parentesi, perché indicano due cose radicalmente diverse:
	\begin{itemize}
		\item Quando si usa la lettera \textbf{senza parentesi}, si intende esclusivamente il \textbf{Poliedro geometrico}, ovvero il puro insieme delle soluzioni ammissibili.
		\item Quando si usa la lettera \textbf{tra parentesi}, ci si riferisce all'\textbf{intero Problema in corso}, che comprende sia lo spazio ammissibile che la Funzione Obiettivo da ottimizzare.
	\end{itemize}
	
	Visivamente e algebricamente, la differenza si riassume in questo modo:
	\[
	\underbrace{P = \{ Ax \leq b \}}_{\text{Il Poliedro}} 
	\qquad \text{vs} \qquad 
	\underbrace{(P) = \begin{cases} \max/\min c^T x \\ x \in P \end{cases}}_{\text{Il Problema completo}}
	\]
	
	In particolar modo, il modello matematico del problema di riferimento era questo:
	\[ 
	\begin{aligned}
		& \begin{cases}
			\max 300x_1 + 200x_2\\
			20x_1 + 10x_2 \leq 950\\
			2x_1 + 3x_2 \leq 120\\
			x_1 + 2x_2 \leq 70\\
			x_{1/2} \geq 0
		\end{cases}
	\end{aligned}
	\]
	
	Le intersezioni che vedete nel grafico (immaginate ci siano dei puntini) si chiamano \textbf{\textit{Vertici}}. Per calcolare l'intersezione tra un vincolo e un altro bisogna mettere a sistema i due o più vincoli in questione e risolverli con l'uguale ($=$). Prendiamo un esempio:
	\[
	\begin{aligned}
		\textbf{Intersezione tra i vincoli 1 e 2}
		\begin{cases}
			20x_1 + 10x_2 = 950\\
			2x_1 + 3x_2 = 120
		\end{cases}
	\end{aligned}
	\]
	Risolvendo questo sistema, si otterranno dei risultati che indicheranno il valore del vertice in cui si intersecano. Ovviamente, i risultati che si ottengono riguardano i valori assunti da $x_1$ e $x_2$ in questo specifico caso di un problema in $R^2$.
	Dopo aver ottenuto i risultati, è fondamentale prendere questi risultati e inserirli nelle equazioni che \textbf{\underline{non}} sono state inserite nel sistema risolvendo però con il $\leq$ e non con l'$=$ come per il sottosistema. Se le disequazioni vengono rispettate, allora il sottosistema di equazioni create rappresenta un \textbf{Vertice} del nostro Poliedro.\\
	
	Quanto è importante il Determinante della matrice di questo sottosistema creato? \textbf{\underline{MOLTISSIMO}}, infatti, se $det(matrice) \neq 0$, allora la soluzione esiste ed è unica e quindi il vertice esiste.\\
	
	Nota bene:
	\begin{itemize}
		\item Ciò che è stato detto a riguardo dei Vertici riguarda solo una grossissima anticipazione di ciò che vedremo più in là. Credo che sia una buona cosa sapere questo in anticipo, ma senza metterci troppa testa perché i dettagli importanti verranno visti dopo.
		\item scrivere "$det(matrice)$" fa parte della sintassi di MATLAB che, ovviamente, durante l'esame si può usare soprattutto per questi dettagli. Semplicemente, crei la matrice A secondo le regole stabilite dal linguaggio, poi si scrive la stringa che vi ho fornito e vi restituisce il risultato.
	\end{itemize}
	
	Come anticipato nella \textit{\underline{Lezione 1}}, nei problemi di PL si esclude l'utilizzo del $<$ e del $>$ dalla realizzazione dei modelli matematici per la seguente definizione:\\\\
	\textit{La soluzione ottima del problema non è dentro al Poliedro, dunque, se esiste, si trova sempre sul bordo}.\\\\
	Dunque, è intuibile che l'esclusione di questi due segni dipende dal fatto che Geometricamente escluderebbero il bordo del poliedro, rendendo impossibile trovare la vera soluzione ottima del problema.
	
	\vspace{1cm}
	
	\textbf{Esercizio di prova}
	\[
	\begin{aligned}
		\begin{cases}
			\max x_1 + x_2\\
			2x_1 + x_2 \leq 7\\
			x_1 + 3x_2 \leq 8\\
			-x_1 \leq 0\\
			-x_2 \leq 0
		\end{cases}
	\end{aligned}
	\]
	Il cui Poliedro è...
	
	\begin{figure}[H]
		\centering
		\begin{tikzpicture}[line cap=round,line join=round]
			\begin{axis}[
				width=12cm, % Gestisce la grandezza fisica sul foglio
				height=12cm,
				axis lines=middle,
				ymajorgrids=true,
				xmajorgrids=true,
				% ZOOM (modifica questi valori per inquadrare diversamente)
				xmin=-5,
				xmax=15,
				ymin=-5,
				ymax=15,
				% Tacche degli assi
				xtick={-5, 0, 5, 10, 15},
				ytick={-5, 0, 5, 10, 15},
				xlabel = {$x_1$},
				ylabel = {$x_2$},
				]
				
				% Riquadro di clip per lo zoom
				\clip(-5,-5) rectangle (15,15);
				
				% VINCOLO 1 (Rosso) 
				% Equivalente a: pspolygon(-2.2,11.4)(-21.72,11.3)(-21.72,-15.72)(11.41,-15.82)(-2.2,11.4)
				\draw[line width=2pt, color=myred, fill=myred, fill opacity=0.15] 
				(-2.2,11.4)--(-21.72,11.3)--(-21.72,-15.72)--(11.41,-15.82)--(-2.2,11.4);
				
				% VINCOLO 2 (Verde)
				% Equivalente a: pspolygon(-21.72,9.90)(-21.72,-15.72)(21.72,-15.72)(21.72,-4.57)(-21.72,9.90)
				\draw[line width=2pt, color=mygreen, fill=mygreen, fill opacity=0.15] 
				(-21.72,9.90)--(-21.72,-15.72)--(21.72,-15.72)--(21.72,-4.57)--(-21.72,9.90);
				
				% VINCOLO DI NON NEGATIVITÀ x1 >= 0 (Grigio)
				\draw[line width=2pt, color=gray, fill=gray, fill opacity=0.1] 
				(0,11.5)--(0,-15.92)--(21.92,-15.92)--(21.92,11.5);
				
				% VINCOLO DI NON NEGATIVITÀ x2 >= 0 (Grigio scuro)
				\draw[line width=2pt, color=darkgray, fill=darkgray, fill opacity=0.1] 
				(-21.92,11.5)--(-21.92,0)--(21.92,0)--(21.92,11.5);
				
			\end{axis}
		\end{tikzpicture}
		\caption{con Regione Ammissibile compresa tra l'asse delle y, l'asse delle x e l'intersezione dei due vincoli. vedi i colori per aiutarti.\\PS: la linea rossa crea quell'angolo in alto per un semplice errore di formattazione grafica.}
	\end{figure}
	
	Per ogni poliedro, esistono le \textit{\textbf{Linee di Isoguadagno}} (o Isocosto). Si ottengono ponendo la FO = K, dunque, in questo caso:\\
	$x_1 + x_2 = k$\\
	Sono delle linee perpendicolari alla direzione del vettore della FO stessa e sono \underline{utilissime} nella ricerca della soluzione ottima. In che modo? Semplice, cerco la linea di Isoguadagno più grande che si trova all'interno del poliedro che, inevitabilmente, sarà sul bordo di quest'ultimo, e quindi seguendo la definizione fornita prima, ossia quella che dice che la soluzione ottima, se esiste, è sul bordo, dimostra che le linee di Isoguadagno servono a trovare proprio l'ottimo del nostro problema.\\
	Una dimostrazione con il grafico:
	\begin{figure}[H]
		\centering
		\begin{tikzpicture}[line cap=round, line join=round]
			\begin{axis}[
				width=12cm,
				height=12cm,
				axis lines=middle,
				ymajorgrids=true,
				xmajorgrids=true,
				% --- IMPOSTAZIONI ZOOM ---
				xmin=-5,
				xmax=15,
				ymin=-5,
				ymax=15,
				% --- GESTIONE DEI NUMERI SUGLI ASSI ---
				xtick={-5, 0, 5, 10, 15},
				ytick={-5, 0, 5, 10, 15},
				xlabel = {$x_1$},
				ylabel = {$x_2$},
				]
				
				% Riquadro di clip per nascondere ciò che esce dallo zoom
				\clip(-5,-5) rectangle (15,15);
				
				% --- AREE AMMISSIBILI E VINCOLI ---
				\draw[line width=2pt, color=myred, fill=myred, fill opacity=0.15]
				(-2.2,11.3)--(-21.72,11.3)--(-21.72,-15.72)--(11.41,-15.72)--(-2.2,11.3);
				
				\draw[line width=2pt, color=mygreen, fill=mygreen, fill opacity=0.15]
				(-21.72,9.91)--(-21.72,-15.72)--(21.72,-15.72)--(21.72,-4.57)--(-21.72,9.91);
				
				% Vincoli di non negatività (Grigio)
				\draw[line width=2pt, color=gray, fill=gray, fill opacity=0.1]
				(0,11.3)--(0,-15.72)--(21.92,-15.72)--(21.92,11.3);
				
				\draw[line width=2pt, color=darkgray, fill=darkgray, fill opacity=0.1]
				(-21.92,11.3)--(-21.92,0)--(21.92,0)--(21.92,11.3);
				
				% --- RETTE (es. Funzioni Obiettivo) ---
				% Sintassi matematica corretta (rimossi i doppi meno "--")
				\draw [line width=2pt, color=orange, domain=-5:15] plot(\x,{1 - \x});
				\draw [line width=2pt, color=orange, domain=-5:15] plot(\x,{0.5 - \x});
				\draw [line width=2pt, color=orange, domain=-5:15] plot(\x,{-\x});
				\draw [line width=2pt, color=orange, domain=-5:15] plot(\x,{-0.5 - \x});
				\draw [line width=2pt, color=orange, domain=-5:15] plot(\x,{1.5 - \x});
				\draw [line width=2pt, color=orange, domain=-5:15] plot(\x,{2 - \x});
				
				% --- ETICHETTE ---
				% Ho spostato le etichette a coordinate visibili (-3) per non farle tagliare fuori dal clip
				\node[font=\scriptsize] at (-3, 6.5) {$eq1$};
				\node[font=\scriptsize] at (-3, 5.5) {$eq2$};
				\node[font=\scriptsize] at (-3, 4.5) {$eq3$};
				\node[font=\scriptsize] at (-3, 3.5) {$eq4$};
				\node[font=\scriptsize] at (-3, 2.5) {$eq5$};
				\node[font=\scriptsize] at (-3, 1.5) {$eq6$};
				
			\end{axis}
		\end{tikzpicture}
		\caption{Le linee di Isoguadagno sono quelle arancioni, perpendicolari alla direzione del vettore della FO e che salgono e/o scendono a seconda del valore K. Dal grafico si può intravedere il fatto che ci sarà un valore K che toccherà i bordi del poliedro. In uno dei punti dei bordi del Poliedro, quindi, c'è l'ottimo del problema. Per questo è molto utile usarle a questo scopo.}
	\end{figure}
	
	Quando parliamo di Poliedri ormai sappiamo che è uno spazio chiuso, sostanzialmente, ma non è sempre così: esistono eccome Poliedri aperti la cui Soluzione Ottima va a $+\infty$. Dunque, ora possiamo anche dire che:\\
	\textit{Se un Poliedro \textbf{non} è chiuso, allora è \underline{illimitato}}.\\
	Un possibile problema che un matematico, un ingegnere e anche banalmente uno studente può trovarsi davanti è la mancanza di Vincoli: si, spesso un Poliedro risulta illimitato perché manca qualche vincolo. Può capitare dato che spesso alcuni sono complicati da trovare nella traccia dell'esercizio e/o in un problema reale. È importante trovarli tutti per proseguire.\\
	
	Allo stesso tempo, esistono anche modelli matematici di PL che, usando GeoGebra o MATLAB, non restituiscono alcun Poliedro. Di sicuro non scompare e di sicuro non hanno sbagliato i due software citati prima dato che sono strumenti molto potenti, piuttosto qui ci troviamo davanti al caso opposto al precedente: se non esiste un Poliedro, sicuramente ci sono troppi vincoli e non esiste una regione ammissibile che rispetta tutti i vincoli del nostro modello. Come per il caso precedente, può capitare di inserire vincoli di troppo o semplicemente il problema non è risolvibile così com'è. Succede anche in situazioni reali.
	
	\vspace{1cm}
	
	Nuovo Problema di PL: \textbf{Problema di Assegnamento}
	
	Ipotizziamo di avere 4 studi di Ingegneria e 4 progetti da realizzare:\\\\
	\begin{center}
		\begin{tabular}{|c|c|c|c|c|}
			\hline
			Progetti $\downarrow$/Studi di ing $\rightarrow$ & \textbf{1} & \textbf{2} & \textbf{3} & \textbf{4}\\
			\hline
			\textbf{1} & 8 & 7 & 9 & 5\\
			\hline
			\textbf{2} & 12 & 13 & 15 & 16\\
			\hline
			\textbf{3} & 9 & 14 & 6 & 10\\
			\hline
			\textbf{4} & 7 & 8 & 11 & 9\\
			\hline
		\end{tabular}
	\end{center}
	Dove i valori interni alla tabella si riferiscono a quanto costerebbe far realizzare un progetto ad un particolare studio. Ad esempio, il progetto 2 costerebbe 16 se lo svolgesse lo studio 4, e così via.
	
	Pappalardo ci fa subito presente una cosa prima della fine della lezione, ossia che il numero possibile di assegnamenti è pari a $\boldsymbol{n!}$ (n fattoriale), mentre invece se inseriamo i vincoli che implicano il fatto che un progetto può esser fatto solo da uno studio e ogni studio può realizzare solo un progetto, gli assegnamenti diventano $\boldsymbol{n^n}$.\\
	
	Fine \textit{Lezione 2}
	
	\separatore\\
	
	\avviso{\textit{\textbf{Questa lezione, come le sue precedenti e successive, può contenere errori e/o imprecisioni. Per il bene di tutti, è buona norma far notare all'autore l'errore o l'imprecisione così da risolvere nel minor tempo possibile. I canali di comunicazione, per eventuali avvisi generali o di errori trovati negli appunti, sono in cima al file. Grazie per la collaborazione e Buono studio!}}}
	
	\section*{Lezione 3 | 23 settembre 2026}
	\subsection*{Approfondimento del problema di assegnamento e studio della Combinazione Convessa e della Combinazione Conica}
	
	Per i prossimi esempi si userà una tabella con valori differenti da quelli della tabella della precedente lezione.
	
	\begin{center}
		\begin{tabular}{|c|c|c|c|c|}
			\hline
			Progetti $\downarrow$/Studi di ing $\rightarrow$ & \textbf{1} & \textbf{2} & \textbf{3} & \textbf{4}\\
			\hline
			\textbf{1} & 8 & 6 & 9 & 4\\
			\hline
			\textbf{2} & 7 & 5 & 6 & 10\\
			\hline
			\textbf{3} & 3 & 8 & 2 & 7\\
			\hline
			\textbf{4} & 4 & 9 & 6 & 5\\
			\hline
		\end{tabular}
	\end{center}
	
	Forse la parte più importante del medesimo problema è stabilire il numero di variabili perché, essendo numerose come vedremo a breve, mancarne una può creare uno sballamento di tutto il modello e del risultato finale.
	
	Innanzitutto, esistono due tipologie di assegnamento:
	\begin{itemize}
		\item \textbf{Cooperativo} (PL): in questo caso, un progetto n potrebbe essere fatto da più di uno studio fino a quando la somma degli studi che collaborano non arriva a 1, come anche uno studio può dedicarsi a più progetti contemporaneamente, cioè: 
		\begin{itemize}
			\item ad esempio, lo studio 1 si occupa dello 0,5 del progetto, lo studio 2 si occupa dello 0,2 e lo studio 3 si occupa dello 0,3. Quindi, $0,5+0,2+0,3=1$, quindi lo studio 4 non collabora nella realizzazione.
			\item Seguendo lo stesso ragionamento, uno studio n può dedicarsi a più di un progetto contemporaneamente (ipotizzando valori e situazioni simili al punto precedente) fino a quando la somma non arriva sempre a 1.
		\end{itemize}
		Come anticipato, qui il numero di assegnamenti possibili è pari a $n!$.
		\item \textbf{Non Cooperativo} (PLI): a differenza del Cooperativo, qui uno studio può dedicarsi ad un solo progetto e ogni progetto deve essere realizzato da un solo studio. Come anticipato, qui il numero di assegnamenti possibili è pari a $n^n$.
	\end{itemize}
	
	Un tipico problema di assegnamento, rappresentato con il nostro tipico Matematichese che Pappalardo ci insegna, è questo:
	\[
	\begin{tabular}{c}
		\text{Assegnamento} \\
		\text{non cooperativo} \\
		\text{in } $\mathbb{R}^{16}$
	\end{tabular}
	\quad
	\begin{cases}
		\begin{aligned} 
			\min \ & 8x_{11} + 6x_{12} + 9x_{13} + 4x_{14} + 7x_{21} + 5x_{22} + 6x_{23} + 10x_{24} \, + \\
			& 3x_{31} + 8x_{32} + 2x_{33} + 7x_{34} + 4x_{41} + 9x_{42} + 6x_{43} + 5x_{44}
		\end{aligned}\\
		x_{11} + x_{12} + x_{13} + x_{14} = 1\\
		x_{21} + x_{22} + x_{23} + x_{24} = 1\\
		x_{31} + x_{32} + x_{33} + x_{34} = 1\\
		x_{41} + x_{42} + x_{43} + x_{44} = 1\\
		x_{11} + x_{21} + x_{31} + x_{41} = 1\\
		x_{12} + x_{22} + x_{32} + x_{42} = 1\\
		x_{13} + x_{23} + x_{33} + x_{43} = 1\\
		x_{14} + x_{24} + x_{34} + x_{44} = 1\\
		x_{ij} \geq 0
	\end{cases}
	\]
	All'esame è caldamente consigliato portarsi dei file (.txt o .md se preferite) che contengono la base di tutti gli script di MATLAB, così dovete solo sostituire i valori del vostro esercizio corrente, dato che scrivere tutto questo porterebbe via una miriade di tempo, soprattutto considerando che il modello matematico "nasconde" una gran parte dei valori che MATLAB dovrà ricevere per trovare l'ottimo.\\
	
	Ecco esattamente come mappare il modello all'interno della funzione \textbf{MATLAB} seguendo alla lettera le regole del \textbf{Formato Primale Standard} (tutto ricondotto a $\leq$ e al $\max$):
	\begin{center}
		\texttt{x = linprog(c, A, b, [], [], [], [])}
	\end{center}
	
	Il nostro problema ha 16 variabili. Dalle 8 uguaglianze originarie ricaveremo 16 vincoli di disuguaglianza (il $\leq$ e il $\geq$ invertito), a cui aggiungeremo i 16 vincoli di non negatività anch'essi invertiti. In totale avremo una matrice \texttt{A} di \textbf{32 righe e 16 colonne}.
	\begin{itemize}
		\item \textbf{\texttt{c}} (Vettore dei costi della F.O. $\rightarrow 1 \times 16$):\\
		Siccome stiamo minimizzando ma \texttt{linprog} cerca nativamente il massimo, dobbiamo invertire i segni della Funzione Obiettivo.\\
		\texttt{c = [-8 -6 -9 -4 -7 -5 -6 -10 -3 -8 -2 -7 -4 -9 -6 -5];}\\
		\textbf{IMPORTANTE PER L'ESAME}: i problemi di assegnamento che ci verranno dati avranno sempre una tabella quadrata perché, è vero che esistono sicuramente problemi con situazioni in cui la tabella non sarebbe quadrata, ma nel nostro corso non vengono trattati, quindi il numero di variabili da inserire nella FO di questo problema sarà sempre $n^2$ con $n$ che si riferisce alla lunghezza di un lato della tabella.
		
		\item \textbf{\texttt{A}} (Matrice dei vincoli $\rightarrow 32 \times 16$):\\
		Popoliamo esplicitamente tutte e 32 le righe. 
		Le prime 8 sono i vincoli $\leq 1$. 
		Le successive 8 sono i vincoli $\geq 1$ invertiti ($\leq -1$). 
		Le ultime 16 sono i vincoli di non negatività invertiti ($-x_{ij} \leq 0$).\\
		\texttt{A = [}\\
		\texttt{1  1  1  1  0  0  0  0  0  0  0  0  0  0  0  0;}\\
		\texttt{0  0  0  0  1  1  1  1  0  0  0  0  0  0  0  0;}\\
		\texttt{0  0  0  0  0  0  0  0  1  1  1  1  0  0  0  0;}\\
		\texttt{0  0  0  0  0  0  0  0  0  0  0  0  1  1  1  1;}\\
		\texttt{1  0  0  0  1  0  0  0  1  0  0  0  1  0  0  0;}\\
		\texttt{0  1  0  0  0  1  0  0  0  1  0  0  0  1  0  0;}\\
		\texttt{0  0  1  0  0  0  1  0  0  0  1  0  0  0  1  0;}\\
		\texttt{0  0  0  1  0  0  0  1  0  0  0  1  0  0  0  1;}\\
		\texttt{-1 -1 -1 -1  0  0  0  0  0  0  0  0  0  0  0  0;}\\
		\texttt{0  0  0  0 -1 -1 -1 -1  0  0  0  0  0  0  0  0;}\\
		\texttt{0  0  0  0  0  0  0  0 -1 -1 -1 -1  0  0  0  0;}\\
		\texttt{0  0  0  0  0  0  0  0  0  0  0  0 -1 -1 -1 -1;}\\
		\texttt{-1  0  0  0 -1  0  0  0 -1  0  0  0 -1  0  0  0;}\\
		\texttt{0 -1  0  0  0 -1  0  0  0 -1  0  0  0 -1  0  0;}\\
		\texttt{0  0 -1  0  0  0 -1  0  0  0 -1  0  0  0 -1  0;}\\
		\texttt{0  0  0 -1  0  0  0 -1  0  0  0 -1  0  0  0 -1;}\\
		\texttt{-1  0  0  0  0  0  0  0  0  0  0  0  0  0  0  0;}\\
		\texttt{0 -1  0  0  0  0  0  0  0  0  0  0  0  0  0  0;}\\
		\texttt{0  0 -1  0  0  0  0  0  0  0  0  0  0  0  0  0;}\\
		\texttt{0  0  0 -1  0  0  0  0  0  0  0  0  0  0  0  0;}\\
		\texttt{0  0  0  0 -1  0  0  0  0  0  0  0  0  0  0  0;}\\
		\texttt{0  0  0  0  0 -1  0  0  0  0  0  0  0  0  0  0;}\\
		\texttt{0  0  0  0  0  0 -1  0  0  0  0  0  0  0  0  0;}\\
		\texttt{0  0  0  0  0  0  0 -1  0  0  0  0  0  0  0  0;}\\
		\texttt{0  0  0  0  0  0  0  0 -1  0  0  0  0  0  0  0;}\\
		\texttt{0  0  0  0  0  0  0  0  0 -1  0  0  0  0  0  0;}\\
		\texttt{0  0  0  0  0  0  0  0  0  0 -1  0  0  0  0  0;}\\
		\texttt{0  0  0  0  0  0  0  0  0  0  0 -1  0  0  0  0;}\\
		\texttt{0  0  0  0  0  0  0  0  0  0  0  0 -1  0  0  0;}\\
		\texttt{0  0  0  0  0  0  0  0  0  0  0  0  0 -1  0  0;}\\
		\texttt{0  0  0  0  0  0  0  0  0  0  0  0  0  0 -1  0;}\\
		\texttt{0  0  0  0  0  0  0  0  0  0  0  0  0  0  0 -1}\\
		\texttt{];}
		
		\item \textbf{\texttt{b}} (Vettore colonna dei termini noti $\rightarrow 32 \times 1$):\\
		Rispetta rigorosamente l'ordine della matrice \texttt{A}: 8 uni positivi, 8 uni negativi e 16 zeri.\\
		\texttt{b = [1; 1; 1; 1; 1; 1; 1; 1; -1; -1; -1; -1; -1; -1; -1; -1; 0; 0; 0; 0; 0; 0; 0; 0; 0; 0; 0; 0; 0; 0; 0; 0];}
		
		\item \textbf{\texttt{Aeq, beq, lb, ub}}:\\
		Poiché abbiamo forzato l'intero modello nel Primale Standard utilizzando unicamente \texttt{A} e \texttt{b}, tutti gli altri parametri obbligatori della sintassi di MATLAB restano vuoti.\\
		\texttt{Aeq = [];}\\
		\texttt{beq = [];}\\
		\texttt{lb = [];}\\
		\texttt{ub = [];}
	\end{itemize}
	
	\textit{Nota pratica}: Avere già pronto sul proprio computer un file con la matrice \texttt{A} schematizzata a blocchi come mostrato qui sopra (magari per problemi tipici fino a $5 \times 5$) fa risparmiare preziosi minuti all'esame, riducendo drasticamente il rischio di perdersi gli zeri per strada.\\
	
	Quello mostrato fin qui implica l'utilizzo dell'assegnamento non cooperativo, e questo dimostrerebbe come sia stato inserito persino il vincolo che obbliga ogni progetto a non essere toccato da più di uno studio. Nella PL si usa l'assegnamento Cooperativo, quindi quello mostrato qui serve principalmente per la PLI, mentre per la PL è uguale togliendo però i vincoli di troppo.\\
	
	Pappalardo ci tiene a darci una definizione importante per l'orale:\\
	\textit{Quando la Regione Ammissibile si allarga, i minimi salgono, quando la Regione Ammissibile si restringe, i minimi scendono}.\\
	
	Ecco alcuni concetti fondamentali per i Poliedri:
	\begin{itemize}
		\item \textbf{\underline{Combinazione Convessa}}:\\
		Dati $k$ punti $x_1, x_2, \dots, x_k \in \mathbb{R}^n$, un generico punto $y \in \mathbb{R}^n$ si definisce \textit{combinazione convessa} di tali punti se:
		\[
		\begin{aligned}
			& \exists \lambda_i \in [0, 1] \quad \text{per } i = 1, \dots, k, \quad \text{con } \sum_{i=1}^k \lambda_i = 1 \quad \text{tale che:} \\
			& y = \sum_{i=1}^k \lambda_i x_i = \lambda_1 x_1 + \lambda_2 x_2 + \dots + \lambda_k x_k
		\end{aligned}
		\]\\
		Da un punto di vista geometrico, l'insieme di tutte le possibili combinazioni convesse di un dato numero di punti rappresenta l'insieme (che in questo caso è a tutti gli effetti un \textbf{Poliedro}) più piccolo esistente in grado di contenere tutti i k punti citati.
		
		\vspace{0.2cm}
		\textbf{Esempio intuitivo:}\\
		Un banale \textbf{segmento} nel piano cartesiano è un poliedro ed è, geometricamente, l'insieme delle combinazioni convesse dei suoi due punti $x_1$ e $x_2$ (situati agli estremi). Ogni singolo punto che giace su quel segmento si ottiene semplicemente variando i valori di $\lambda_1$ e $\lambda_2$.
		
		\begin{figure}[H]
			\centering
			\begin{tikzpicture}[scale=1.5]
				% Coordinate dei 2 punti estremi
				\coordinate (x1) at (0, 0);
				\coordinate (x2) at (4, 1.5);
				
				% Un generico punto sul segmento (es. con lambda_1 = 0.5 e lambda_2 = 0.5)
				\coordinate (y) at (2, 0.75); 
				
				% Disegno il segmento (il Poliedro unidimensionale)
				\draw[thick, color=black] (x1) -- (x2);
				
				% Disegno i pallini per i punti e le etichette
				\filldraw (x1) circle (1pt) node[below left] {$x_1$};
				\filldraw (x2) circle (1pt) node[above right] {$x_2$};
				
				% Disegno il punto interno generato dalla combinazione convessa
				\filldraw[color=blue] (y) circle (1pt) node[below right, text=black] {$y$};
				
			\end{tikzpicture}
			\caption{Il segmento è l'insieme di tutte le Combinazioni Convesse dei suoi estremi $x_1$ e $x_2$. Un generico punto $y$ si ottiene variando i pesi $\lambda_1$ e $\lambda_2$.}
		\end{figure}
		
		\vspace{0.2cm}
		\textbf{Altro esempio (Punto interno):}\\
		Supponiamo di avere 4 punti nel piano, di cui tre formano i vertici di un triangolo e il quarto ($x_4$) si trova posizionato all'interno di essi. L'insieme di tutte le loro combinazioni convesse genererà proprio quel triangolo. Anche se $x_4$ non "partecipa" alla creazione del bordo (non è un vertice), il triangolo risultante è comunque la Combinazione Convessa dell'intero set, perché rappresenta fedelmente il Poliedro più piccolo in assoluto in grado di contenerli tutti e 4.
		
		\begin{figure}[H]
			\centering
			\begin{tikzpicture}[scale=1.5]
				% Coordinate dei 4 punti
				\coordinate (x1) at (0, 0);
				\coordinate (x2) at (2, 2);
				\coordinate (x3) at (4, -0.5);
				\coordinate (x4) at (2.2, 0.4); % Punto interno
				
				% Disegno il perimetro del triangolo (il Poliedro)
				\draw[thick, color=black] (x1) -- (x2) -- (x3) -- cycle;
				
				% Riempimento leggero per far capire che è un'area (opzionale, decommenta se ti piace)
				% \fill[gray, opacity=0.1] (x1) -- (x2) -- (x3) -- cycle;
				
				% Disegno i pallini per i punti e le etichette
				\filldraw (x1) circle (1pt) node[below left] {$x_1$};
				\filldraw (x2) circle (1pt) node[above right] {$x_2$};
				\filldraw (x3) circle (1pt) node[below right] {$x_3$};
				\filldraw (x4) circle (1pt) node[right] {$x_4$};
			\end{tikzpicture}
			\caption{Combinazione Convessa di 4 punti dove $x_4$ è interno al poliedro generato.}
		\end{figure}
		
		\textbf{\underline{Attenzione a non confondersi!}}\\
		Prendendo come riferimento i 4 punti dell'esempio appena visto, cosa succederebbe se decidessimo di "chiudere" il perimetro collegando i punti in questo modo: $x_1 \rightarrow x_2 \rightarrow x_4 \rightarrow x_3 \rightarrow x_1$? 
		
		Otterremmo una figura con una "rientranza". Questa figura \textbf{NON} rappresenta la Combinazione Convessa dei 4 punti. 
		
		Il motivo geometrico e matematico è rigoroso: se disegnassimo i bordi in quel modo, tutta l'area compresa tra il punto $x_4$ e il segmento diretto $x_2-x_3$ verrebbe esclusa dal nostro insieme. Ma un insieme, per essere \textit{convesso} (e quindi per essere una vera Combinazione Convessa), deve rispettare la regola d'oro: \textbf{presi due punti qualsiasi appartenenti all'insieme, l'intero segmento che li unisce deve essere contenuto all'interno dell'insieme stesso}. 
		
		Come si vede chiaramente nel grafico sottostante, se provassimo a unire $x_2$ e $x_3$ con un segmento, questo finirebbe "fuori" dal nostro perimetro rientrante. Dunque, non è una Combinazione Convessa valida.
		
		\begin{figure}[H]
			\centering
			\begin{tikzpicture}[scale=1.5]
				% Coordinate dei 4 punti (le stesse di prima)
				\coordinate (x1) at (0, 0);
				\coordinate (x2) at (2, 2);
				\coordinate (x3) at (4, -0.5);
				\coordinate (x4) at (2.2, 0.4);
				
				% Disegno il perimetro ERRATO (rientrante, non convesso)
				\draw[thick, color=red] (x1) -- (x2) -- (x4) -- (x3) -- cycle;
				
				% Disegno il segmento che dimostra che non è convesso
				\draw[thick, dashed, color=black] (x2) -- (x3);
				
				% Etichetta per l'area esclusa erroneamente (spostata all'esterno con anchor=west)
				\node[color=red, align=left, anchor=west] at (3.3, 1.4) {\footnotesize Segmento\\[-1mm]\footnotesize fuori dall'insieme\\[-1mm]\footnotesize (Area esclusa!)};
				
				% Disegno i pallini per i punti e le etichette
				\filldraw (x1) circle (1pt) node[below left] {$x_1$};
				\filldraw (x2) circle (1pt) node[above right] {$x_2$};
				\filldraw (x3) circle (1pt) node[below right] {$x_3$};
				\filldraw (x4) circle (1pt) node[below left] {$x_4$}; % Etichetta spostata per non incrociare le linee
			\end{tikzpicture}
			\caption{Esempio di figura NON convessa. La "rientranza" in $x_4$ esclude un'area, violando la regola fondamentale della convessità: il segmento $x_2-x_3$ passa fuori dall'insieme.}
		\end{figure}
		\item \textbf{\underline{Combinazione Conica}}:\\
		Dati $k$ punti $x_1, x_2, \dots, x_k \in \mathbb{R}^n$, un generico punto $y \in \mathbb{R}^n$ si definisce \textit{combinazione conica} di tali punti se:
		\[
		\begin{aligned}
			& \exists \lambda_i \geq 0 \quad \text{per } i = 1, \dots, k, \quad \text{tale che:} \\
			& y = \sum_{i=1}^k \lambda_i x_i = \lambda_1 x_1 + \lambda_2 x_2 + \dots + \lambda_k x_k
		\end{aligned}
		\]
		
		\vspace{0.2cm}
		\textbf{Differenze chiave rispetto alla Convessa:}
		\begin{itemize}
			\item I pesi $\lambda_i$ non sono più limitati superiormente a 1, ma possono assumere qualsiasi valore reale purché non negativo ($\lambda_i \geq 0$).
			\item Cade il vincolo fondamentale $\sum \lambda_i = 1$. Non è più richiesta alcuna somma fissa.
		\end{itemize}
		
		\vspace{0.3cm}
		Da un punto di vista geometrico, l'insieme di tutte le possibili combinazioni coniche di un dato insieme di vettori genera un \textbf{cono convesso} avente il vertice nell'origine. 
		
		Mentre la combinazione convessa crea un "recinto" chiuso e limitato tra i punti stessi (come un segmento o un triangolo), la combinazione conica si "espande" all'infinito partendo dall'origine e passando attraverso e oltre quei punti, proprio come il fascio di luce di una torcia.
		
		\vspace{0.2cm}
		\textbf{Rappresentazione Grafica (Il Cono):}\\
		Nel grafico sottostante vediamo la differenza visiva fondamentale della combinazione conica. Partendo dall'origine degli assi, i vettori diretti verso i punti $x_1$ e $x_2$ non si fermano, ma si estendono all'infinito creando un "fascio" (l'area ombreggiata). 
		Tutta questa regione infinita è a tutti gli effetti un \textbf{Poliedro illimitato}.
		
		Come evidenziato dall'esempio pratico, i punti $x_1$ e $x_2$ appartengono ovviamente all'insieme da loro stessi generato (sono combinazioni coniche valide, indicate con $\checkmark$), mentre un generico punto $x_3$, trovandosi al di fuori di questo fascio, \textbf{non} può in alcun modo essere generato da una combinazione conica di $x_1$ e $x_2$ (dunque ne resta escluso, $\times$).
		
		\begin{figure}[H]
			\centering
			\begin{tikzpicture}[scale=1.5]
				% Assi cartesiani
				\draw[->, thin, gray] (-0.5, 0) -- (4.5, 0) node[right, text=black] {asse $x$};
				\draw[->, thin, gray] (0, -0.5) -- (0, 3) node[above, text=black] {asse $y$};
				
				% Coordinate dei punti
				\coordinate (O) at (0,0);
				\coordinate (x1) at (1.2, 2.2);
				\coordinate (x2) at (2.5, 0.8);
				\coordinate (x3) at (3.2, 0.3); % Il punto che sta "fuori"
				
				% Raggi infiniti (moltiplicati per farli estendere oltre i punti)
				\coordinate (r1) at (1.5, 2.75);
				\coordinate (r2) at (3.8, 1.21);
				
				% Area della Combinazione Conica (Poliedro illimitato)
				% Uso un colore grigio/azzurro leggero per simulare il tratteggio del quaderno
				\fill[blue, opacity=0.08] (O) -- (r1) -- (r2) -- cycle;
				
				% Vettori (frecce) dall'origine ai punti generatori
				\draw[->, thick, color=black] (O) -- (x1);
				\draw[->, thick, color=black] (O) -- (x2);
				
				% Estensione dei raggi tratteggiata per far capire che vanno all'infinito
				\draw[dashed, thick, color=black] (x1) -- (r1);
				\draw[dashed, thick, color=black] (x2) -- (r2);
				
				% Pallini per i punti ed etichette (le spunte e le X come nel quaderno)
				\filldraw (x1) circle (1pt) node[left] {$x_1 (\checkmark)$};
				\filldraw (x2) circle (1pt) node[below] {$x_2 (\checkmark)$};
				\filldraw[color=red] (x3) circle (1pt) node[right, text=red] {$x_3 (\times \text{ fuori})$};
				
				% Etichetta Poliedro con la freccia, simile all'appunto a mano
				\draw[<-, thick, color=darkgray] (1.8, 1.4) -- (2.8, 2.2) node[right, align=left] {\textbf{Poliedro illimitato}\\[-1mm]\footnotesize (perché delimitato\\[-1mm]\footnotesize da semispazi ma \footnotesize non è chiuso)};
				
			\end{tikzpicture}
			\caption{Combinazione Conica generata dai punti $x_1$ e $x_2$. L'area ombreggiata rappresenta un poliedro illimitato. Il punto $x_3$ è esterno al cono, pertanto non è combinazione conica di $x_1$ e $x_2$.}
		\end{figure}
	\end{itemize}
	
	\vspace{0.8cm}
	\separatore
	\vspace{0.3cm}
	
	\subsection*{Teorema di Rappresentazione di Weyl}
	
	Avendo definito le combinazioni convesse e coniche, possiamo introdurre uno dei teoremi più importanti della Ricerca Operativa che unisce i due concetti.
	
	Dato un poliedro $P$ (che sappiamo essere definito dalla matrice $A$ e dal vettore colonna $b$ nel modello $Ax \leq b$), il teorema afferma che esistono sempre:
	\begin{itemize}
		\item Un insieme di punti (vertici) $V = \{v_1, v_2, \dots, v_k\} \in P$.
		\item Un insieme di direzioni (raggi) $E = \{e_1, e_2, \dots, e_p\}$, dove $p$ è semplicemente l'indice massimo dell'insieme $E$ e non c'entra nulla con la lettera $P$ del Poliedro.
	\end{itemize}
	
	Allora l'intero Poliedro $P$ si può esprimere come la \textbf{somma della combinazione convessa dell'insieme $V$ e della combinazione conica dell'insieme $E$}.
	
	$P=Conv(V)+Conp(E)$
	
	In "Matematichese" si scrive così:
	\[
	\begin{aligned}
		P = \left\{ x \in \mathbb{R}^n \;\middle|\; x = \underbrace{\sum_{i=1}^k \lambda_i v_i}_{\text{Comb. Convessa di } V} + \underbrace{\sum_{j=1}^p \mu_j e_j}_{\text{Comb. Conica di } E} \right\}
	\end{aligned}
	\]
	Rispettando ovviamente i vincoli tipici che abbiamo studiato per le due combinazioni:
	\[
	\begin{aligned}
		& \lambda_i \geq 0 \quad \forall i = 1, \dots, k \quad \text{e} \quad \sum_{i=1}^k \lambda_i = 1 \\
		& \mu_j \geq 0 \quad \forall j = 1, \dots, p
	\end{aligned}
	\]
	
	\vspace{0.3cm}
	\textbf{Significato della somma di insiemi:}\\
	Come si può intuire dalla formula, la combinazione convessa esprime \textit{tutte} le possibili combinazioni convesse dei punti di $V$, mentre la combinazione conica esprime \textit{tutte} le possibili combinazioni coniche delle direzioni di $E$. 
	
	In matematica, la somma di due insiemi di questo tipo (nota come \textit{Somma di Minkowski}) rappresenta esattamente \textbf{l'insieme di tutte le possibili somme che si possono fare prendendo un elemento dal primo insieme e un elemento dal secondo}. 
	
	Geometricamente, significa che prendiamo la figura "chiusa" generata dalla combinazione convessa dei vertici e la facciamo "scivolare" (traslare) lungo le direzioni infinite generate dalla combinazione conica, ottenendo così la forma esatta e completa del nostro poliedro originario.
	
	\vspace{0.3cm}
	\textbf{Cosa significa concretamente tutto questo?}\\
	In maniera geometrica e visiva, sommare questi due insiemi significa prendere \textbf{ogni singolo punto} dell'insieme $V$ (generato dalla combinazione convessa) e trattarlo esattamente come se fosse una "nuova origine". Da ognuna di queste nuove origini di partenza viene fatto scaturire l'intero fascio infinito di direzioni creato dall'insieme $E$ (la combinazione conica). 
	
	È letteralmente come se prendessimo il cono di direzioni (generato da $E$) e ne facessimo coincidere il vertice di partenza con ogni singolo punto valido dell'insieme $V$: l'unione di tutte queste proiezioni "spazza" l'intero spazio generando la forma esatta e completa del nostro Poliedro originario. Questo è il significato concreto della frase detta prima, ossia "\textit{Allora l'intero Poliedro $P$ si può esprimere come la \textbf{somma della combinazione convessa dell'insieme $V$ e della combinazione conica dell'insieme $E$}.}"
	\begin{center}
		\textit{Nota: Questo teorema e la relativa costruzione geometrica verranno ripresi e approfonditi nel dettaglio con degli esempi pratici nella prossima lezione.}
	\end{center}
	
	Fine \textit{Lezione 3}
	
	\separatore\\
	
	\avviso{\textit{\textbf{Questa lezione, come le sue precedenti e successive, può contenere errori e/o imprecisioni. Per il bene di tutti, è buona norma far notare all'autore l'errore o l'imprecisione così da risolvere nel minor tempo possibile. I canali di comunicazione, per eventuali avvisi generali o di errori trovati negli appunti, sono in cima al file. Grazie per la collaborazione e Buono studio!}}}
	
	\section*{Lezione 4 | 24 settembre 2026}
	\subsection*{Comprensione del vero utilizzo del teorema di Weyl e Teorema fondamentale della PL}
	
	Il Teorema di Rappresentazione di Weyl non è solo una nozione teorica, ma lo strumento pratico su cui si fonda gran parte della Ricerca Operativa. La sua vera utilità è permetterci di \textbf{passare costantemente e senza perdita di informazioni} dalla fredda rappresentazione algebrica (quella "da calcolatore", definita dalla matrice $A$ e dal vettore colonna $b$) alla rappresentazione basata sui vertici e sulle direzioni.
	
	\vspace{0.2cm}
	\textbf{Esempio pratico sulla Somma di Insiemi:}\\
	Riprendendo il concetto finale della scorsa lezione, cerchiamo di visualizzare cosa significa sommare due insiemi (la cosiddetta \textit{Somma di Minkowski}). 
	Ipotizziamo che il nostro insieme dei vertici $V$ sia composto solo da due punti, $v_1$ e $v_2$. La loro combinazione convessa ($Conv(V)$) genera un banale segmento limitato. 
	Ora ipotizziamo che l'insieme delle direzioni $E$ contenga un unico raggio $e_1$ che punta verso l'alto. La sua combinazione conica ($Cone(E)$) è una retta infinita che parte dall'origine. 
	
	Cosa succede applicando il Teorema di Weyl ($P = Conv(V) + Cone(E)$)? 
	Stiamo prendendo ogni singolo punto del segmento limitato e lo stiamo usando come origine per far scaturire la direzione infinita verso l'alto. Geometricamente, è come se "trascinassimo" l'intero segmento verso l'infinito superiore. Il risultato finale è un nuovo poliedro illimitato a forma di "nastro" infinito.
	
	\vspace{0.4cm}
	\textbf{Il Triangolo delle Rappresentazioni}\\
	L'obiettivo primario di questo corso è rendervi capaci di muovervi agilmente tra queste tre forme. Molto spesso, un problema che risulta incomprensibile o troppo calcoloso nella sua Forma Algebrica, può essere risolto in pochissimi passaggi se tradotto e analizzato nella sua Forma Geometrica.
	
	\begin{figure}[H]
		\centering
		\begin{tikzpicture}[
			>=stealth,
			block/.style={rectangle, draw=titlecolor, fill=titlecolor!5, 
				text width=5cm, text centered, rounded corners=4pt, 
				minimum height=1.8cm, line width=1pt, font=\small},
			arrow/.style={<->, line width=1.2pt}
			]
			
			% Nodi (I tre blocchi) posizionati con coordinate assolute per evitare sovrapposizioni
			\node[block] (algebrica) at (0, 0) {\textbf{Forma Algebrica}\\[2mm] Esempio:\\ $Ax \leq b, \quad x \geq 0$};
			
			\node[block] (weyl) at (-4.8, -4.5) {\textbf{Forma di Weyl}\\[2mm] Esempio:\\ $P = Conv(V) + Cone(E)$};
			
			\node[block] (geometrica) at (4.8, -4.5) {\textbf{Forma Geometrica}\\[2mm] Risoluzione del problema guardandolo\\ visivamente e graficamente.};
			
			% Collegamenti (Frecce bi-direzionali senza testo)
			\draw[arrow, color=myblue] (algebrica) -- (weyl);
			
			\draw[arrow, color=myred] (algebrica) -- (geometrica);
			
			\draw[arrow, color=mygreen] (weyl) -- (geometrica);
			
		\end{tikzpicture}
	\end{figure}
	
	Nota bene:
	\begin{itemize}
		\item All'esame orale Pappalardo non ti chiederà di rappresentare graficamente un problema in $\mathbb{R}^6$, ma neanche in $\mathbb{R}^3$, dato che è impossibile farlo a mano. Dunque, è utile sicuramente ragionare in $\mathbb{R}^2$, all'orale una domanda riguardante questo argomento non è improbabile, anzi.
		\item All'esame scritto è importante però conoscere bene la logica di passaggio da una forma all'altra perché lui comunque vuole sempre tutti i passaggi. Meglio sapere almeno di cosa si sta parlando.
	\end{itemize}
	
	\vspace{0.4cm}

	\textbf{Il trucco pratico sull'illimitatezza del Poliedro:}\\
	Possiamo visualizzare concretamente questo fenomeno applicando la Funzione Obiettivo ($c^T x$) all'intera espressione del Teorema di Weyl. Sostituendo la $x$ con la somma degli insiemi, otteniamo:
	\[
	c^T x = c^T \left( \sum_{i=1}^k \lambda_i v_i + \sum_{j=1}^p \mu_j e_j \right)
	\]
	
	\vspace{0.2cm}
	Distribuendo il vettore dei costi $c^T$ all'interno delle due sommatorie, la separazione tra la parte limitata e la parte potenzialmente illimitata del problema diventa evidente:
	\[
	c^T x = \underbrace{\sum_{i=1}^k \lambda_i (c^T v_i)}_{\text{Parte limitata (Vertici)}} + \underbrace{\sum_{j=1}^p \mu_j (c^T e_j)}_{\text{Parte divergente (Direzioni)}}
	\]
	Il professore ci ha fornito una chiave di lettura immediata e un trucco pratico per analizzare la natura del Poliedro e del problema di ottimizzazione semplicemente guardando gli elementi di questa formula:
	\begin{itemize}
		\item \textbf{Poliedro limitato:} Se l'insieme delle direzioni $E$ è vuoto ($E = \emptyset$), l'intera seconda sommatoria (quella relativa alla combinazione conica) scompare. Il Poliedro risulta formato unicamente dalla combinazione convessa dei suoi vertici ($P = Conv(V)$). Questo ci dà la certezza matematica assoluta che il poliedro è chiuso e \textbf{limitato}.
		\item \textbf{Poliedro illimitato e divergenza:} La presenza stessa di una direzione $e_j$ rende il poliedro strutturalmente \textbf{illimitato}. Il trucco sta nell'analizzare il prodotto tra la Funzione Obiettivo (il vettore $c^T$) e le direzioni della combinazione conica. Se esiste almeno una direzione $e_j$ tale per cui il prodotto scalare è maggiore di zero:
		\[
		\begin{aligned}
			\sum_{j=1}^p \mu_j e_j \qquad \text{con:} \quad c^T e_j > 0
		\end{aligned}
		\]
		allora l'intero membro non si annulla. Questo ha un significato concreto fortissimo: significa che muovendoci all'infinito lungo quel raggio, il valore della nostra funzione obiettivo continuerà a crescere all'infinito. È l'indicazione precisa che non solo il poliedro è illimitato, ma che il problema \textbf{non ammette un ottimo finito} (la soluzione diverge a $+\infty$).
	\end{itemize}
	
	\vspace{0.3cm}
	\textbf{La massimizzazione algebrica della Funzione Obiettivo:}\\
	Come evidenziato dal professore, questa trasformazione (ovvero tradurre il $c^T x$ generico nei prodotti mirati $c^T v_i$ e $c^T e_j$) non è un semplice vezzo teorico, ma rappresenta il \textbf{vero e proprio metodo algebrico per massimizzare la Funzione Obiettivo}. 
	
	Invece di dover cercare la soluzione ottima "alla cieca" tra gli infiniti punti $x$ del Poliedro, questa scomposizione ci permette di controllare il massimo in maniera sistematica:
	\begin{itemize}
		\item \textbf{Controllo delle direzioni:} Se anche per un solo raggio il prodotto $c^T e_j$ risulta strettamente maggiore di zero, sappiamo subito che all'aumentare di $\mu_j$ il valore schizzerà a $+\infty$. Abbiamo trovato la divergenza senza fare calcoli complessi.
		\item \textbf{Controllo dei vertici:} Se invece nessuna direzione ci fa "esplodere" la funzione (quindi tutti i $c^T e_j \leq 0$), la componente di destra si azzera ai fini della massimizzazione. A quel punto, per trovare il massimo assoluto dell'intero problema, ci basterà calcolare i vari $c^T v_i$ e prendere banalmente il vertice che restituisce il valore più alto, come vedremo in uno degli esempi.
	\end{itemize}
	
	\vspace{0.4cm}
	\textbf{Dimostrazione dell'ottimo sui vertici:}\\
	Ipotizziamo ora di trovarci in un caso in cui il problema non diverga all'infinito. La componente delle direzioni si annulla (ovvero $\sum \mu_j e_j = 0$ ai fini del calcolo) e ogni punto $x$ del Poliedro può essere espresso come sola combinazione convessa dei vertici:
	\[
	x = \sum_{i=1}^k \lambda_i v_i
	\]
	
	Applicando la Funzione Obiettivo, cerchiamo il massimo di $c^T x$:
	\[
	\max c^T x = \max \left( \sum_{i=1}^k \lambda_i (c^T v_i) \right)
	\]
	
	Sviluppando la sommatoria, l'espressione diventa banalmente:
	\[
	\lambda_1(c^T v_1) + \lambda_2(c^T v_2) + \dots + \lambda_k(c^T v_k)
	\]
	
	Sappiamo per definizione che i pesi $\lambda_i$ sono tutti non negativi e che la loro somma deve fare rigorosamente 1 ($\sum \lambda_i = 1$). Grazie a questa proprietà, possiamo applicare un trucco algebrico fondamentale (una \textit{maggiorazione}): una combinazione convessa di più valori non potrà \textbf{mai} superare il valore del termine più grande moltiplicato per la somma dei pesi. 
	
	Se chiamiamo $v_{max}$ il vertice che restituisce il valore più alto in assoluto tra tutti i $v_i$, possiamo sostituirlo agli altri e "raccoglierlo":
	\[
	\sum_{i=1}^k \lambda_i (c^T v_i) \leq (c^T v_{max}) \underbrace{(\lambda_1 + \lambda_2 + \dots + \lambda_k)}_{= 1} = c^T v_{max}
	\]
	
	Poiché $v_{max}$ è esso stesso un punto appartenente al Poliedro $P$ (è uno dei suoi vertici), abbiamo appena dimostrato matematicamente che il valore massimo ottenibile sull'intero poliedro è uguale al valore massimo ottenibile sui soli vertici:
	\[
	\max_{x \in P} c^T x = \max_{v_i \in V} c^T v_i
	\]
	
	Questa uguaglianza è la prova definitiva dell'equivalenza tra il \textbf{Formato Primale Standard} e la \textbf{Rappresentazione di Weyl}. Ci conferma la regola d'oro: cercare l'ottimo sull'intero dominio del poliedro equivale a cercarlo esplorando unicamente i suoi vertici. 
	
	\vspace{0.2cm}
	\textit{Nota: In presenza di una componente conica con direzioni di crescita illimitata ($\sum \mu_j e_j$ con $c^T e_j > 0$), la medesima struttura logica porterebbe banalmente la funzione a scalare verso $+\infty$, confermando l'illimitatezza del Poliedro $P$ e del problema stesso.}
	
	\vspace{0.3cm}
	
	\subsection*{Esercizi utili come esempio}
	Di seguito analizzeremo alcuni casi studio pratici per dimostrare il legame diretto tra la Geometria, l'Algebra (matrici $A$ e $b$) e la Rappresentazione di Weyl (insiemi $V$ ed $E$).
	
	\begin{enumerate}[label=\textbf{Esercizio \arabic*:}, itemsep=15pt, parsep=5pt, leftmargin=*]
		
		\item \textbf{Il Quadrato Unitario (Poliedro Limitato)}\\
		Abbiamo un problema definito dalla matrice $A$ e dal vettore colonna $b$:
		\[
		A = \begin{bmatrix} 
			-1 & 0 \\ 
			0 & -1 \\ 
			1 & 0 \\ 
			0 & 1 
		\end{bmatrix}, \qquad 
		b = \begin{bmatrix} 
			0 \\ 
			0 \\ 
			1 \\ 
			1 
		\end{bmatrix}
		\]
		Traduciamo l'espressione algebrica $Ax \leq b$ in vincoli geometrici veri e propri, ottenendo un sistema lineare:
		\[
		\begin{cases}
			x_1 \geq 0 \\
			x_2 \geq 0 \\
			x_1 \leq 1 \\
			x_2 \leq 1
		\end{cases}
		\]
		
		Essendo un poliedro perfettamente limitato, l'insieme delle direzioni è vuoto ($E = \emptyset$). I suoi vertici sono banalmente $V = \{v_1, v_2, v_3, v_4\}$ con $v_1=(0,0)$, $v_2=(1,0)$, $v_3=(1,1)$ e $v_4=(0,1)$. L'intero poliedro è generato esclusivamente dalla Combinazione Convessa di questi 4 punti ($P = Conv(V)$).
		
		\begin{figure}[H]
			\centering
			\begin{tikzpicture}
				\begin{axis}[
					width=8cm, height=8cm, axis lines=middle, ymajorgrids=true, xmajorgrids=true,
					xmin=-0.5, xmax=1.8, ymin=-0.5, ymax=1.8, xtick={0,1}, ytick={0,1},
					xlabel={$x_1$}, ylabel={$x_2$}
					]
					\draw[line width=2pt, color=myblue, fill=myblue, fill opacity=0.2] (0,0) -- (1,0) -- (1,1) -- (0,1) -- cycle;
					
					\filldraw (0,0) circle (2pt) node[below left] {$v_1 (0,0)$};
					\filldraw (1,0) circle (2pt) node[below right] {$v_2 (1,0)$};
					\filldraw (1,1) circle (2pt) node[above right] {$v_3 (1,1)$};
					\filldraw (0,1) circle (2pt) node[above left] {$v_4 (0,1)$};
					
					\node[align=center, fill=white, draw=titlecolor, rounded corners=3pt, inner sep=4pt] at (0.5, 0.5) {\footnotesize $P = Conv(V)$};
				\end{axis}
			\end{tikzpicture}
		\end{figure}
		
		\item \textbf{Il Quarto Quadrante (Poliedro Illimitato)}\\
		In questo caso la forma algebrica è molto compatta. Abbiamo la matrice $A$ e il vettore $b$:
		\[
		A = \begin{bmatrix} 
			-1 & 0 \\ 
			0 & 1 
		\end{bmatrix}, \qquad 
		b = \begin{bmatrix} 
			0 \\ 
			0 
		\end{bmatrix}
		\]
		Traduciamo l'espressione $Ax \leq b$ nei vincoli geometrici, ottenendo il sistema:
		\[
		\begin{cases}
			x_1 \geq 0 \\
			x_2 \leq 0
		\end{cases}
		\]
		Siamo nel quadrante in basso a destra del piano cartesiano. Essendo il poliedro illimitato, la Rappresentazione di Weyl richiederà l'insieme delle direzioni $E$. 
		L'unico vertice disponibile per partire è l'origine, quindi $V = \{v_1\}$ con $v_1=(0,0)$. Da lì scaturiscono due direzioni all'infinito: $e_1=(1,0)$ (lungo l'asse $x_1$ positivo) e $e_2=(0,-1)$ (lungo l'asse $x_2$ negativo).
		
		\begin{figure}[H]
			\centering
			\begin{tikzpicture}
				\begin{axis}[
					width=8cm, height=8cm, axis lines=middle, ymajorgrids=true, xmajorgrids=true,
					xmin=-1, xmax=4, ymin=-4, ymax=1, xtick={0,2,4}, ytick={0,-2,-4},
					xlabel={$x_1$}, ylabel={$x_2$}
					]
					\draw[line width=0pt, fill=myblue, fill opacity=0.15] (0,0) -- (4,0) -- (4,-4) -- (0,-4) -- cycle;
					
					\draw[->, line width=2pt, color=myred] (0,0) -- (3,0) node[above] {$e_1 (1,0)$};
					\draw[->, line width=2pt, color=myred] (0,0) -- (0,-3) node[left] {$e_2 (0,-1)$};
					
					\filldraw (0,0) circle (2pt) node[above left] {$v_1 (0,0)$};
				\end{axis}
			\end{tikzpicture}
		\end{figure}
		
		\item \textbf{Il Cono Traslato}\\
		Cosa succede se il vertice di partenza non è l'origine? Consideriamo $V = \{v_1\}$ con $v_1=(1,2)$ e direzioni $E = \{e_1, e_2\}$ con $e_1=(1,1)$ e $e_2=(1,0)$. 
		Come studiato nel Teorema di Weyl, la combinazione conica usa \textbf{ogni punto dell'insieme $V$ come nuova origine}. In questo caso, le frecce delle direzioni non partono da $(0,0)$, ma si agganciano al vertice $v_1$, "trascinando" il cono nel piano.
		
		\begin{figure}[H]
			\centering
			\begin{tikzpicture}
				\begin{axis}[
					width=8cm, height=8cm, axis lines=middle, ymajorgrids=true, xmajorgrids=true,
					xmin=-0.5, xmax=5, ymin=-0.5, ymax=5, xtick={0,1,2,3,4}, ytick={0,1,2,3,4},
					xlabel={$x_1$}, ylabel={$x_2$}
					]
					\draw[line width=0pt, fill=myblue, fill opacity=0.15] (1,2) -- (4,5) -- (5,5) -- (5,2) -- cycle;
					
					\draw[->, line width=2pt, color=myred] (1,2) -- (3,4) node[above left] {$e_1 (1,1)$};
					\draw[->, line width=2pt, color=myred] (1,2) -- (4,2) node[below] {$e_2 (1,0)$};
					
					\filldraw (1,2) circle (2pt) node[above left] {$v_1 (1,2)$};
				\end{axis}
			\end{tikzpicture}
		\end{figure}
		
		\item \textbf{La Direzione Ridondante}\\
		Dato $V=\{v_1\}$ con $v_1=(0,0)$ e direzioni $e_1=(1,1)$ ed $e_2=(1,0)$. Ci viene fornita una terza direzione $e_3=(1,0.5)$ e ci viene chiesto se serve a generare il poliedro. 
		La risposta è \textbf{NO}. Dal grafico è palese che il raggio $e_3$ transita internamente allo spazio già delimitato e generato da $e_1$ ed $e_2$. Matematicamente significa che $e_3$ è già una Combinazione Conica di $e_1$ ed $e_2$, risulta quindi \textit{ridondante} e può essere scartata dall'insieme $E$.
		
		\begin{figure}[H]
			\centering
			\begin{tikzpicture}
				\begin{axis}[
					width=8cm, height=8cm, axis lines=middle, ymajorgrids=true, xmajorgrids=true,
					xmin=-0.5, xmax=3, ymin=-0.5, ymax=3, xtick={0,1,2}, ytick={0,1,2},
					xlabel={$x_1$}, ylabel={$x_2$}
					]
					\draw[line width=0pt, fill=mygreen, fill opacity=0.15] (0,0) -- (3,3) -- (3,0) -- cycle;
					
					\draw[->, line width=2pt, color=myblue] (0,0) -- (2,2) node[above left] {$e_1 (1,1)$};
					\draw[->, line width=2pt, color=myblue] (0,0) -- (2,0) node[below] {$e_2 (1,0)$};
					
					\draw[->, line width=2pt, dashed, color=myred] (0,0) -- (2,1) node[right, font=\footnotesize] {$e_3$ (Ridondante)};
					
					\filldraw (0,0) circle (2pt) node[above left] {$v_1 (0,0)$};
				\end{axis}
			\end{tikzpicture}
		\end{figure}
		
		\item \textbf{Il Segmento come Poliedro Limitato}\\
		Dato l'insieme $V=\{(1,1), (2,2)\}$ ed $E = \emptyset$. L'obiettivo qui è distruggere il pregiudizio visivo che un poliedro debba per forza essere una figura con un'area (come un quadrato o un triangolo). 
		La sola combinazione convessa di due punti genera un segmento. Un segmento ha un perimetro chiuso dai suoi stessi estremi, pertanto è a tutti gli effetti un Poliedro chiuso e limitato.
		
		\begin{figure}[H]
			\centering
			\begin{tikzpicture}
				\begin{axis}[
					width=8cm, height=8cm, axis lines=middle, ymajorgrids=true, xmajorgrids=true,
					xmin=-0.5, xmax=3.5, ymin=-0.5, ymax=3.5, xtick={0,1,2,3}, ytick={0,1,2,3},
					xlabel={$x_1$}, ylabel={$x_2$}
					]
					% Disegnato come un vero segmento senza frecce (non è un vettore)
					\draw[line width=2.5pt, color=myblue] (1,1) -- (2,2);
					
					\filldraw (1,1) circle (2.5pt) node[below right] {$v_1 (1,1)$};
					\filldraw (2,2) circle (2.5pt) node[above left] {$v_2 (2,2)$};
				\end{axis}
			\end{tikzpicture}
		\end{figure}
		
		\item \textbf{Il Tubo Infinito (Somma Minkowski Avanzata)}\\
		Questo esercizio è uno "scherzo logico" che inganna l'occhio. Dati $V = \{v_1, v_2\}$ con $v_1=(0,1)$ e $v_2=(0,2)$. Ci viene chiesto di aggiungere direzioni per creare un poliedro illimitato da $-\infty$ a $+\infty$ sull'asse $x$. 
		
		L'errore comune è pensare che $E$ abbia direzioni "alle altezze $y=1$ e $y=2$". In realtà, per definizione, \textbf{le direzioni dell'insieme $E$ sono semplici vettori ancorati all'origine}. L'insieme $E$ contiene semplicemente le direzioni destra e sinistra: $e_1=(1,0)$ ed $e_2=(-1,0)$. 
		
		La vera magia avviene applicando Weyl (la \textit{Somma di Minkowski}): le direzioni $e_1$ ed $e_2$ vengono duplicate e fatte scaturire \textbf{da tutti i punti} della combinazione convessa dei vertici (ossia l'intero segmento tra $v_1$ e $v_2$). Trascinando il segmento a destra e a sinistra verso l'infinito, si crea visivamente questo "tubo" o nastro orizzontale.
		
		\begin{figure}[H]
			\centering
			\begin{tikzpicture}
				\begin{axis}[
					width=12cm, height=6cm, axis lines=middle, ymajorgrids=true, xmajorgrids=true,
					xmin=-5, xmax=5, ymin=-0.5, ymax=3, xtick={-4,-2,0,2,4}, ytick={0,1,2},
					xlabel={$x_1$}, ylabel={$x_2$}
					]
					% L'area infinita del tubo
					\draw[line width=0pt, fill=orange, fill opacity=0.2] (-5,1) rectangle (5,2);
					
					% Il segmento dei vertici Conv(V)
					\draw[line width=3pt, color=myblue] (0,1) -- (0,2);
					
					% I vertici
					\filldraw (0,1) circle (2.5pt) node[below right, text=black] {$v_1 (0,1)$};
					\filldraw (0,2) circle (2.5pt) node[above right, text=black] {$v_2 (0,2)$};
					
					% Le direzioni applicate al centro per dare l'idea del trascinamento
					\draw[->, line width=2pt, color=myred] (0,1.5) -- (3,1.5) node[above] {$e_1 (1,0)$};
					\draw[->, line width=2pt, color=myred] (0,1.5) -- (-3,1.5) node[above] {$e_2 (-1,0)$};
				\end{axis}
			\end{tikzpicture}
		\end{figure}
		
		\item \textbf{Verifica Analitica e Calcolo dell'Ottimo}\\
		In quest'ultimo caso non ci viene fornita la matrice $A$, ma passiamo direttamente all'approccio analitico avendo a disposizione gli insiemi della Rappresentazione di Weyl e il vettore della Funzione Obiettivo $c$:
		\[
		V = \left\{ v_1=\begin{pmatrix} 1 \\ 1 \end{pmatrix}, v_2=\begin{pmatrix} 2 \\ 3 \end{pmatrix}, v_3=\begin{pmatrix} -1 \\ 0 \end{pmatrix} \right\}, \quad 
		E = \left\{ e_1=\begin{pmatrix} 0 \\ 1 \end{pmatrix}, e_2=\begin{pmatrix} 1 \\ 3 \end{pmatrix} \right\}, \quad 
		c = \begin{bmatrix} 4 \\ -3 \end{bmatrix}
		\]
		Il professore ribadisce un concetto fondamentale: l'insieme dei vertici è sempre contenuto nel poliedro ($V \subset P$). 
		
		Per trovare il massimo, non basta calcolare subito i vertici. Dobbiamo applicare rigorosamente la regola studiata precedentemente procedendo in due step.
		
		\textbf{Step 1: Controllo della divergenza sulle direzioni ($E$)}\\
		Verifichiamo innanzitutto se il problema esplode a $+\infty$ calcolando i prodotti scalari $c^T e_j$:
		\begin{itemize}
			\item $c^T e_1 = (4 \cdot 0) + (-3 \cdot 1) = -3 \leq 0$
			\item $c^T e_2 = (4 \cdot 1) + (-3 \cdot 3) = 4 - 9 = -5 \leq 0$
		\end{itemize}
		Dato che nessuna direzione produce un valore strettamente positivo, siamo certi che \textbf{il problema non diverge}. L'ottimo finito esiste e si trova sicuramente su uno dei vertici!
		
		\textbf{Step 2: Calcolo dell'ottimo sui vertici ($V$)}\\
		Solo ora siamo autorizzati a calcolare il prodotto $c^T v_i$ per ogni vertice e prendere il valore massimo:
		\begin{itemize}
			\item $c^T v_1 = (4 \cdot 1) + (-3 \cdot 1) = 4 - 3 = 1$ \quad $\leftarrow \textbf{MASSIMO}$
			\item $c^T v_2 = (4 \cdot 2) + (-3 \cdot 3) = 8 - 9 = -1$
			\item $c^T v_3 = (4 \cdot -1) + (-3 \cdot 0) = -4 - 0 = -4$
		\end{itemize}
		La soluzione ottima del problema è quindi il valore $1$, che si ottiene in corrispondenza del vertice $v_1 = (1,1)$.
		
		\begin{figure}[H]
			\centering
			\begin{tikzpicture}
				\begin{axis}[
					width=10cm, height=9cm, axis lines=middle, ymajorgrids=true, xmajorgrids=true,
					xmin=-2, xmax=4, ymin=-1, ymax=6, xtick={-2,-1,0,1,2,3,4}, ytick={-1,0,1,2,3,4,5,6},
					xlabel={$x_1$}, ylabel={$x_2$}
					]
					% Area del Poliedro illimitato generato dalla Somma di Minkowski
					\draw[line width=0pt, fill=myblue, fill opacity=0.1] (-1,6) -- (-1,0) -- (1,1) -- (2,3) -- (3,6) -- cycle;
					
					% Conv(V) - Il triangolo generato dalla Combinazione Convessa dei vertici
					\draw[line width=1.5pt, color=myblue, dashed, fill=myblue, fill opacity=0.25] (-1,0) -- (1,1) -- (2,3) -- cycle;
					
					% Vertici
					\filldraw[color=mygreen] (1,1) circle (3pt) node[below right, text=black] {$v_1 (1,1)$ \textbf{(Ottimo)}};
					\filldraw (2,3) circle (2pt) node[right] {$v_2 (2,3)$};
					\filldraw (-1,0) circle (2pt) node[above left] {$v_3 (-1,0)$};
					
					% Direzioni E applicate ai vertici estremi per mostrare l'illimitatezza
					\draw[->, line width=2pt, color=myred] (-1,0) -- (-1,3) node[left] {$e_1 (0,1)$};
					\draw[->, line width=2pt, color=myred] (2,3) -- (3,6) node[right] {$e_2 (1,3)$};
					
				\end{axis}
			\end{tikzpicture}
			\caption{Somma di Minkowski dell'Esercizio 7. Il triangolo tratteggiato è $Conv(V)$, mentre le frecce rosse sono le direzioni $Cone(E)$ applicate ai vertici estremi. Il punto $v_1$ massimizza la Funzione Obiettivo. All'inizio può sembrare strano, ma sapendo che il vettore $c$ va praticamente verso il basso, oltre che algebricamente, si capisce come il punto $v_1$ sia l'ottimo a discapito del punto $v_2$ che poteva sembrarlo a primo impatto. È sempre importante vedere dove va la Funzione Obiettivo.}
		\end{figure}
		\vspace{0.4cm}
		\textbf{Nota Geometrica sulla Ridondanza Visiva:}\\
		Osservando il grafico dell'Esercizio 7, si nota che le direzioni dell'insieme $E$ non sono state tracciate a partire dal vertice $v_1(1,1)$, e sui vertici $v_2$ e $v_3$ è stata applicata una sola direzione per ciascuno. 
		
		Questo avviene per un principio di \textbf{ridondanza geometrica}. Sebbene la Somma di Minkowski imponga matematicamente di applicare l'intero cono delle direzioni a \textit{ogni} punto di $Conv(V)$, ai fini della rappresentazione grafica ci interessa tracciare unicamente l'involucro esterno (i bordi) del Poliedro. 
		
		\begin{itemize}
			\item \textbf{Sul vertice interno ($v_1$):} Qualsiasi direzione fatta scaturire da qui punterebbe direttamente verso l'interno dell'area già evidenziata, risultando visivamente superflua.
			\item \textbf{Sui vertici esterni ($v_2$ e $v_3$):} Tracciamo esclusivamente il raggio necessario a costruire la "parete" esterna del poliedro (ad esempio, solo $e_1$ su $v_3$). L'altra direzione finirebbe per transitare all'interno del poliedro, venendo "assorbita" dallo spazio già delimitato.
		\end{itemize}
		
		Questa semplice intuizione ci permette di tradurre la complessa regola algebrica di Weyl in un disegno rapido, pulito e privo di informazioni sovrapposte.
		
	\end{enumerate}
	
	\subsection*{Teorema Fondamentale della Programmazione Lineare}
	
	Tutto ciò che abbiamo dimostrato fino a questo momento, a partire dalle combinazioni convesse fino alla scomposizione della Funzione Obiettivo, converge nel risultato teorico in assoluto più importante di questa prima parte del corso. 
	
	Dato un problema di Programmazione Lineare rigorosamente in \textbf{Formato Primale Standard}, e definiti $V$ (insieme dei vertici) ed $E$ (insieme delle direzioni) come gli insiemi della rappresentazione del suo Poliedro $P$, il teorema afferma che \textbf{si verificherà sempre e soltanto uno dei seguenti tre casi}:
	
	\begin{itemize}[label=\textcolor{titlecolor}{\Large\textbullet}, itemsep=8pt, parsep=2pt, leftmargin=1.5cm]
		
		\item \textbf{Caso 1: Poliedro Vuoto (Inammissibilità)}\\
		Il problema non ammette alcuna soluzione ammissibile perché il Poliedro è banalmente vuoto ($P = \emptyset$). Nella pratica geometrica e analitica, questo accade tipicamente quando ci sono \textit{troppi vincoli} o vincoli severamente contraddittori tra loro, rendendo impossibile l'esistenza di uno spazio geometrico condiviso che li rispetti tutti.
		
		\item \textbf{Caso 2: Problema Illimitato (Divergenza)}\\
		Il problema ammette soluzioni ammissibili, ma il valore della Funzione Obiettivo diverge verso $+\infty$. Come dimostrato precedentemente, questo si verifica quando il Poliedro è aperto lungo direzioni favorevoli alla Funzione Obiettivo. Spesso, nei problemi reali, questo è il sintomo della dimenticanza di alcuni limiti fisici, traducendosi in \textit{troppo pochi vincoli} a "chiudere" la regione ammissibile.
		
		\item \textbf{Caso 3: Esistenza dell'Ottimo Finito (Soluzione sui vertici)}\\
		Se non si verificano i due casi patologici precedenti (ovvero il Poliedro esiste e la funzione non diverge), allora il problema ammette rigorosamente un ottimo finito. In questo scenario, il teorema garantisce matematicamente che \textbf{uno dei punti dell'insieme $V$ (dunque uno dei vertici del Poliedro) è sicuramente la soluzione ottima} del nostro problema.
		
	\end{itemize}
	\vspace{0.4cm}
	\textbf{Condizioni algebriche di verifica:}\\
	A livello puramente operativo, il professore ha sintetizzato il riconoscimento di questi tre scenari in regole matematiche immediate:
	\begin{itemize}
		\item Ci troviamo nel \textbf{Caso 1} se l'insieme dei vertici è vuoto ($V = \emptyset$).
		\item Ci troviamo nel \textbf{Caso 2} se, calcolando i prodotti tra la Funzione Obiettivo e le direzioni, esiste anche solo un risultato strettamente positivo ($c^T e_j > 0$).
		\item Siamo nel \textbf{Caso 3} quando non ricadiamo nei due scenari precedenti. In questa situazione, per trovare la soluzione ottima è sufficiente calcolare i prodotti $c^T v_i$ per tutti i vertici dell'insieme $V$ e selezionare il valore massimo.
	\end{itemize}
	
\end{document}